\chapter{Motivations}
\section{Motivation and approach for HyDrop}
At the start of our project in 2018, single-cell sequencing methods had become established research tools, but were still expensive and complicated enough to preclude routine use. At the time, academic scRNA-seq protocols were either low in throughput\sjcitep[,]{islamCharacterizationSinglecellTranscriptional2011, ramskoldFulllengthMRNASeqSinglecell2012, hashimshonycelseqsinglecellrnaseq2012, sasagawa2013quartz} or low in sensitivity\sjcitep[,]{macoskohighlyparallelgenomewide2015, gierahnseqwellportablelowcost2017, caoComprehensiveSinglecellTranscriptional2017,kleinDropletBarcodingSinglecell2015} and reserved for use by specialised labs (see \autoref{sect:single_cell_sequencing_tech}). The majority of cells were sequenced using the 10x Chromium, a commercial device/kit which packaged inDrop's droplet microfluidic cell isolation strategy with aspects of Drop-seq and SMART-seq's RNA library preparation into a user-friendly package\sjcitep[.]{zhengmassivelyparalleldigital2017} The Chromium provided a convenient - and more importantly \textit{reliable} - way to generate high-quality single-cell RNA-seq libraries\sjcitep[]{mereuBenchmarkingSinglecellRNAsequencing2020} and struck a comfortable balance between throughput and sensitivity.

We made three observations:
\begin{enumerate}[noitemsep]
	\item At a cost of ~\euro{\num{10000}} per 8 samples of 10 000 cells, the 10x scRNA kit was ill-suited for exploratory research.
	\item The kit operated as a black box, with limited possibility for modification.
	\item The kit's sensitivity was around 10 times higher than inDrop and Drop-seq's, implying that the sensitivity ceiling for droplet-based methods had not yet been reached by open-source methods.
\end{enumerate}

In addition, a droplet-based single-cell ATAC sequencing method was not available yet, again constraining researchers to either low-sensitivity\sjcitep[]{cusanovichMultiplexSinglecellProfiling2015} or low-throughput\sjcitep[]{buenrostrosinglecellchromatinaccessibility2015,chenrapidrobustmethod2018} methods. Since droplet methods had provided a balanced trade-off between throughput and sensitivity that fit our lab's scRNA-seq needs\sjcitep[,]{davie2018single} we set out to develop an equivalent for scATAC-seq. We also understood that combined scRNA/ATAC-seq would be necessary to deconvolve the individual and (cross-) combined effects of gene expression and chromatin accessibility in future experiments. A better understanding of how to design and build a droplet-based scATAC-seq method would be a good stepping stone to building multiomics technologies. Specifically, we wanted to build an open-source droplet platform capable of both scRNA-seq and scATAC-seq, at higher sensitivities than inDrop and Drop-seq. This work had already been initiated before my arrival at the lab, and I spent my first year as a master thesis student contributing to this project. Further researching and developing this method, which we later named HyDrop, occupied the first half of this PhD project.

\section{Motivation and approach for scATAC-seq benchmark}
Throughout the development of HyDrop-ATAC, we were often confronted with questions surrounding the method's sensitivity (the number of fragments detected per cell) and specificity (the enrichment of these fragments in enhancers and near transcription start sites, collectively forming chromatin accessibility "peaks"). Such questions arose both internally and externally, and echoed previous calls for an in-depth method performance exploration since the release of Bio-Rad\supershortnotecite{lareaudropletbasedcombinatorialindexing2019} and 10x Genomics'\supershortnotecite{satpathyMassivelyParallelSinglecell2019} scATAC-seq products in 2019. At the time, a number of open-source scATAC-seq protocols existed\sjcitep[,]{cusanovichsinglecellatlasvivo2018,chenrapidrobustmethod2018}, but required significant levels of expertise (e.g. automated liquid handling for scaling) or non-commercially available components (e.g. indexed Tn5). Both BioRad and 10x's scATAC-seq solutions were more user-friendly and handled the most technically challenging aspect, microfluidic droplet generation, on benchtop devices (10x's Chromium Controller and Bio-Rad's ddSEQ). However, the Chromium controller already enjoyed a wide presence in large research centres due to the success of 10x's scRNA-seq products, and the Fluidigm C1 had fallen out of use due to throughput and per-cell price considerations. As a result, 10x dominated the commercial scATAC-seq market from 2016 onward, leading to \textasciitilde{1000} publications compared to ddSEQ's \textasciitilde{70} by 2025.

We reasoned that an exact quantification of the performance differences among commercial and non-commercial protocols could stimulate informed decision making. Several such benchmarks had been performed for scRNA-seq methods: the earliest by Wu and Quake\sjcitep[,]{wu2013quantitative} one by Ziegenhain and Enard\sjcitep[,]{ziegenhainComparativeAnalysisSingleCell2017} and a larger, more recent study by Mereu and Heyn\sjcitep[.]{mereuBenchmarkingSinglecellRNAsequencing2020} In the Mereu benchmark, a reference sample was distributed to collaborating labs, which then each generated at least one scRNA-seq dataset. This reference consisted of a mixture of PBMCs, mouse colon cells, HEK293T, 3T3 and dog MDCK cells. This approach complicated genome alignment and further post-processing, and we therefore opted to limit the reference sample to a mixture of human PBMC extracted from one male and one female donor. The use of only liquid tissue also simplified sample handling: PBMC can be distributed frozen and thawed slowly before nuclei extraction. Finally, a large reference dataset of single-cell PBMC chromatin accessibility did not yet exist, ensuring that our benchmark's final dataset could also provide value as a public resource.
